\documentclass[10pt,a4paper]{amsart}
\usepackage[margin=19mm]{geometry}
\usepackage{amsmath,amssymb,mathtools}
\usepackage[scr=rsfs,cal=euler]{mathalfa}
\usepackage{xfrac}
\usepackage[unicode,hidelinks,bookmarks=false]{hyperref}
\newcommand{\ie}{i.e.\ }
\newcommand{\cf}{cf.\ }
\newcommand{\resp}{respectively\ }
\newcommand{\Subsection}{\subsection}
\providecommand{\nobreakdash}{\nobreak\hbox{-}}
\setlength{\emergencystretch}{2em}
\multlinegap0pt
\usepackage{bm}
\usepackage{bbm}
\usepackage{dsfont}
\usepackage{xspace}
\usepackage{cancel}
\usepackage{mathtools}
\usepackage{amsthm}
\makeatletter
\@ifundefined{lemma}{\newtheorem{lemma}{Lemma}}{}
\makeatother
\newcommand{\ZZ}{\mathbb{Z}}
\newcommand{\cD}{\mathcal{D}}
\newcommand{\cP}{\mathcal{P}}
\title[Sum of consecutive powers as a perfect power]{Sum of consecutive powers as a perfect power}
\author{Angelos Koutsianas, Nikos Tzanakis}
\date{Source: arXiv:2605.18348v1 (CC BY 4.0)}
\begin{document}
\maketitle

\noindent\emph{Source and licence.} Sum of consecutive powers as a perfect power. Version arXiv:2605.18348v1 (CC BY 4.0), distributed under the Creative Commons Attribution 4.0 International licence, as recorded in the arXiv licence metadata for this exact version and shown on the abstract page https://arxiv.org/abs/2605.18348v1. Full rights notice: licenses/arxiv-2605.18348-CC-BY-4.0.txt.\par\smallskip

\noindent\textit{Editorial scope.} Counted unit: the Lemma whose source label is \texttt{lem:equations\_system} in the article (source0.tex lines 312--320, source label \texttt{lem:equations\_system}), delivered together with the article's own complete original proof of it (source0.tex lines 322--386, one \texttt{proof} environment of 65 lines). No mathematics is written, repaired or re-derived: statement and proof are the article's own text line for line. Required context retained uncounted: eq:main\_symmetric (lines 166--169); lem:quadratic\_factor (lines 274--284); eq:fk (lines 280--283). Every definition, display and label of those lines is retained unchanged. Omitted from this package and not redistributed: 1--165, 170--273, 284--311, 321--321, 387--1097. Nothing inside a retained excerpt is omitted or altered: every retained body is delivered exactly as the article's lines are written, so no omission receipt of any kind accompanies this package. Citations: the retained excerpts carry no citation command at all, so no bibliography entry of the article is transcribed and no reference is redistributed. Reference adaptation: none was needed and none was made; every reference inside the delivered unit resolves inside it, so no unresolved reference or citation is produced. Printed slips kept literal: no printed word, symbol, hypothesis or number of the counted statement or of its proof was altered. Layout and class: the article's own class is \texttt{amsart} with packages including fullpage, tabu, amssymb, amsmath, amscd, amsbsy, comment, enumerate, xy, hyperref, mathrsfs, color, mathtools, caption; the wrapper is the lane's pinned \texttt{amsart} editorial wrapper at 10pt on A4, which restates the theorem-like environments the retained text needs and adds the article's own macro definitions that the retained text needs (\texttt{\textbackslash ZZ}, \texttt{\textbackslash cD}, \texttt{\textbackslash cP}), with extra wrapper packages (bm, bbm, dsfont, xspace, cancel, mathtools). The delivered numbering is the unit's own. Assets: no retained span contains a figure environment, a graphics-inclusion command, a PDF-page inclusion, a table or a TikZ picture; no image file is copied into this package, the package declares no asset dependency and no image file is redistributed with it. Licence: version arXiv:2605.18348v1 of this article is distributed under the Creative Commons Attribution 4.0 International licence, as recorded in the arXiv licence metadata for this exact version and shown on the abstract page https://arxiv.org/abs/2605.18348v1. A full rights notice is delivered at licenses/arxiv-2605.18348-CC-BY-4.0.txt. Characters: every retained excerpt is pure ASCII, so the delivered engine, XeLaTeX with its default Latin Modern text font, reports no missing character.
\begin{equation}\label{eq:main_symmetric}
   \frac{(x-1)^k + (x+1)^k}{2} = 2^{k-1}y^n, \quad 
   x,y\ge 3\text{ odd},\quad n\ge 3,
\end{equation}

\begin{lemma}\label{lem:quadratic_factor}
We have $g_k(t)\in\ZZ[t]$ and 
\begin{equation*}
    g_k(t) = (t^2 + 1)f_k(t^2),
\end{equation*}
where $f_k(t)\in\ZZ[t]$ and $\deg(f)=\frac{k}{2} - 1$. In particular,
\begin{equation}\label{eq:fk}
    f_k(t) = \sum_{i=0}^{(k/2-1)/2}2^{k/2-(2i+1)}
    \binom{k/2}{2i+1}(t+1)^{2i}t^{(k/2-(2i+1))/2}.
\end{equation}
\end{lemma}

\begin{lemma}\label{lem:equations_system}
Let $k\ge 6$ and suppose $(x,y,n)$ is a solution of \eqref{eq:main_symmetric} with $n>2v_p(k)$ for every prime $p\equiv 1\pmod 4$ dividing $k$. Let $\cP$ be the set of primes which divide $k/2$ and are congruent to $1\pmod 4$, if any such primes exist; otherwise set $\cP=\emptyset$. Define $\cD$ to be the set consisting of $1$ and the positive divisors of $k/2$ whose prime divisors belong to $\cP$ (if $\cP\ne\emptyset$). Then
\begin{align}
    x^2 + 1 & = 2\frac{d_2}{d_1}y_1^n,     \label{eq:express_x^2+1}
    \\
    f_k(x^2) & = 2^{k-2}\frac{d_1}{d_2}y_2^n,   \label{eq:express_f_k(x^2)}
\end{align}
where $d_1,d_2\in\cD$ are relatively prime, $y_1,y_2$ are relatively prime positive odd integers and $y_1y_2=y$.
\end{lemma}

\begin{proof}
By \eqref{eq:main_symmetric} and  Lemma \ref{lem:quadratic_factor} we have $(x^2+1)f_k(x^2)=2^{k-1}y^n$. Since $x$ is odd, $v_2(x^2+1)=1$, therefore $(x^2+1,f_k(x^2)):=2d$ with $d$ odd. If $d>1$, every prime $p$ dividing $d$ also divides $x^2+1$, therefore $p\equiv 1\pmod 4$. Also,
\begin{equation} \label{eq:eq-system-1}
    \frac{x^2+1}{2d}\cdot\frac{f_k(x^2)}{2d}=2^{k-3}\frac{y^n}{d^2}.
\end{equation}
From \eqref{eq:fk} with $t=x^2$ we get $2d\mid 2^{k/2-1}(k/2)$, hence $d\mid (k/2)$; we have already seen that, if $d>1$, every prime divisor of $d$ is $\equiv 1\pmod 4$, therefore $d\in \cD$. Moreover, $v_p(d)\le v_p(k)<n/2$ for every odd prime $p$, a fact that will be used below.

If $d=1$, then we put $d_1=d_2=1$ and \eqref{eq:eq-system-1} 
immediately implies the relations \eqref{eq:express_x^2+1} and
\eqref{eq:express_f_k(x^2)} with appropriate $y_1,y_2$ as in
the announcement of the lemma.

We consider now the case $d>1$.
From the last displayed equation it is clear that $y^n/d^2\in\ZZ$. 
Let $\cP_0\neq\emptyset$ be the set of primes dividing $d$ 
($\cP_0\subseteq\cP$), and write 
$d=\prod_{p\in \cP_0}p^{e_p}$ for its canonical factorization. 
Then $y=y'\prod_{p\in \cP_0}p^{m_p}$  where $m_p:=v_{p}(y)>0$ and  
$(y',\prod_{p\in \cP_0}p)=1$. Consequently, by \eqref{eq:eq-system-1},
%
\begin{equation} \label{eq:eq-system-2}
    \frac{x^2+1}{2d}\cdot\frac{f_k(x^2)}{2d}= 
 2^{k-3}y'^n\prod_{p\in \cP_0}p^{nm_p-2e_p}
=2^{k-3}\prod_{p\in \cP_0}p^{-2e_p}\cdot
\left(y'\prod_{p\in \cP_0}p^{m_p}\right)^n,
\end{equation}
%
Note that, for every $p\in \cP_0$, the $p$-exponent on the right-hand side is positive because $m_p>0$ and $n>2v_p(d)=2e_p$. On the other hand, the two factors on the left-hand side are relatively prime, therefore, for every $p\in \cP_0$, precisely one of them has positive $p$-exponent. It follows that we can set  $\cP_0=\cP_1\sqcup \cP_2$, where $\cP_1$ is the set of those primes $p\in \cP_0$ with non-zero $p$-exponent in $(x^2+1)/(2d)$ and, analogously, $\cP_2$ is the set of the primes $p\in \cP_0$ whose $p$-exponent in $f_k(x^2)/(2d)$ is non-zero. Note that $\cP_1$ or $\cP_2$ (but not both) may be empty. Now, from \eqref{eq:eq-system-2} it follows that 

\begin{equation}  \label{eq:express_the_2_factors}
\frac{x^2+1}{2d}=
\prod_{p\in \cP_1}p^{-2e_p}\cdot
\left(y'_1\prod_{p\in \cP_1}p^{m_p}\right)^n,
\quad
 \frac{f_k(x^2)}{2d} = 2^{k-3}
    \prod_{p\in \cP_2}p^{-2e_p}\cdot
    \left(y'_2\prod_{p\in \cP_2}p^{m_p}\right)^n
\end{equation}
%
with $(y'_1,y'_2)=1$, $y'_1$ odd and $y'_1y'_2=y'$.
Putting
%
\begin{equation*}
    y_i = y'_i\prod_{p\in \cP_i}p^{m_p},\quad i=1,2,
\end{equation*}
%
we see that $y_1$ is odd and
$y_1y_2=y'_1y'_2\prod_{p\in \cP_0}p^{m_p}=y'\prod_{p\in \cP_0}p^{m_p}=y$.
Then, by \eqref{eq:express_the_2_factors},
\begin{equation*}
x^2+1=2d\prod_{p\in\cP_1}p^{-2e_p}y_1^n,\quad
f_k(x^2)=2^{k-3}d\prod_{p\in\cP_2}p^{-2e_p}y_1^n,
\end{equation*}
and, since $d=\prod_{p\in\cP_1\sqcup\cP_2}p^{e_p}$, it is 
straightforward to see that
\begin{equation*}
x^2+1=2\prod_{p\in\cP_1}p^{-e_p}\prod_{p\in\cP_2}p^{e_p}y_1^n,
\quad
f_k(x^2)= 2^{k-3}\prod_{p\in\cP_1}p^{e_p}\prod_{p\in\cP_2}p^{-e_p}y_2^n.
\end{equation*}
We complete the proof by setting 
\begin{equation*}
d_1= \prod_{p\in\cP_1}p^{e_p},\quad d_2 = \prod_{p\in\cP_2}p^{e_p}.
\end{equation*}
\end{proof}

\end{document}
